\section{Теоретические основы алгоритмов сортировки}

\subsection{Пузырьковая сортировка}

Пузырьковая сортировка~--- простейший алгоритм, который последовательно
сравнивает соседние элементы и меняет их местами, если они стоят в
неправильном порядке~\cite{Knuth}. Несмотря на свою простоту, алгоритм
обладает плохой временной сложностью и практически не применяется в
промышленных системах.

Количество сравнений в худшем случае для массива из \(n\) элементов:

\begin{equation}
    \label{eq:bubble}
    C_{\text{bubble}}(n) = \sum_{i=1}^{n-1} i = \frac{n(n-1)}{2} = O(n^2).
\end{equation}

\subsection{Сортировка слиянием}

Сортировка слиянием~--- алгоритм типа «разделяй и властвуй»
(\textit{divide and conquer})~\cite{Cormen}. Массив рекурсивно
разбивается на две половины, каждая сортируется отдельно, затем
отсортированные половины сливаются в один массив.

Рекуррентное соотношение для числа операций:

\begin{equation}
    \label{eq:mergesort}
    T(n) = 2T\!\left(\frac{n}{2}\right) + \Theta(n),
\end{equation}

решением которого по теореме Мастера является \(T(n) = \Theta(n \log n)\).

\subsection{Быстрая сортировка}

Быстрая сортировка (\textit{quicksort}) также относится к классу
«разделяй и властвуй». Выбирается опорный элемент (\textit{pivot}),
массив разбивается на два подмассива~--- меньше и больше опорного~---
и каждый сортируется рекурсивно~\cite{Hoare}.

Ожидаемое число сравнений при случайном выборе опорного элемента:

\begin{equation}
    \label{eq:quicksort}
    E[C(n)] = 2n \ln n - \frac{2(n-1)}{n} \approx 1{,}386\, n \log_2 n.
\end{equation}

\subsection{Сравнительная таблица}

В таблице~\ref{tab:complexity} приведено сравнение временной и
пространственной сложности рассматриваемых алгоритмов.

\begin{table}[H]
    \centering
    \caption{Сравнение сложности алгоритмов сортировки}
    \label{tab:complexity}
    \begin{tabular}{|l|c|c|c|c|}
        \hline
        \textbf{Алгоритм} &
        \textbf{Лучший случай} &
        \textbf{Средний случай} &
        \textbf{Худший случай} &
        \textbf{Память} \\
        \hline
        Пузырьковая & $O(n)$        & $O(n^2)$      & $O(n^2)$      & $O(1)$        \\
        \hline
        Слияние     & $O(n \log n)$ & $O(n \log n)$ & $O(n \log n)$ & $O(n)$        \\
        \hline
        Быстрая     & $O(n \log n)$ & $O(n \log n)$ & $O(n^2)$      & $O(\log n)$   \\
        \hline
    \end{tabular}
\end{table}
